% Current synthesis for Lectures 5--7. Included by synthesis.tex.

% ============================================================
\section{Lecture 5 --- Self-Financing, Replication, and Black--Scholes}
% ============================================================

\lecturesummary{It\^o's formula identifies the dynamics of a candidate
price. Self-financing specifies how a trading strategy earns its gains.
Matching those gains with a European claim gives the Black--Scholes PDE
and its delta hedge. The Gaussian solution produces the call formula and
its Greeks. The basic model in this lecture has no dividends.}

\subsection{From It\^o's formula to a pricing equation}

For a stock with historical dynamics
\[
\dd S_t=\mu S_t\dd t+\sigma S_t\dd W_t^{\Pbb},
\qquad
\dd B_t=rB_t\dd t,\qquad B_0=1,
\]
It\^o's formula applied directly to \(V_t=v(t,S_t)\) gives
\[
\dd V_t
=\left(v_t+\mu S_tv_s+\tfrac12\sigma^2S_t^2v_{ss}\right)\dd t
+\sigma S_tv_s\dd W_t^{\Pbb}.
\]
The second derivative term comes from \(\dd[S]_t=\sigma^2S_t^2\dd t\).
For path-dependent payoffs, a running integral such as
\(A_t=\int_0^tS_u\dd u\) can be added to the state. This idea returns in
Lecture~7.

\subsection{The self-financing budget}

Let \(\Delta_t\) be stock shares and \(\beta_t\) bank-account units. The
portfolio value is \(X_t=\Delta_tS_t+\beta_tB_t\). If holdings change at
a trading date \(t_i\), the new shares must be paid for from the bank
account:
\[
S_{t_i}(\Delta_i-\Delta_{i-1})
+B_{t_i}(\beta_i-\beta_{i-1})=0.
\]
The holdings chosen before the next price move earn that move. Passing
from these left-point gains to continuous time gives the
self-financing condition
\[
\boxed{
X_t=X_0+\int_0^t\Delta_u\dd S_u+\int_0^t\beta_u\dd B_u,
\qquad
\dd X_t=\Delta_t\dd S_t+r(X_t-\Delta_tS_t)\dd t.}
\]
The value identity \(X=\Delta S+\beta B\) by itself does not impose
self-financing. For continuous finite-variation holdings, the product rule
also gives \(S_t\dd\Delta_t+B_t\dd\beta_t=0\). For more general
semimartingale holdings, a covariation term must be included in that
rebalancing identity; the stochastic-integral gain equation remains the
definition.

\subsection{Replication, the PDE, and verification}

Suppose a European claim pays \(g(S_T)\) and has smooth candidate price
\(v(t,S_t)\), with \(v(T,s)=g(s)\). A replicating portfolio with
\(X_t=v(t,S_t)\) has gain
\[
\dd X_t
=\left[rv+(\mu-r)S_t\Delta_t\right]\dd t
+\sigma S_t\Delta_t\dd W_t^{\Pbb}.
\]
Matching its Brownian coefficient with that of \(\dd v(t,S_t)\) gives
\(\Delta_t=v_s(t,S_t)\). Matching drifts then cancels \(\mu\) and yields
\[
\boxed{
v_t+rs\,v_s+\tfrac12\sigma^2s^2v_{ss}-rv=0,
\qquad v(T,s)=g(s).}
\]
The cash holding is
\[
\boxed{\beta_t=\frac{v(t,S_t)-S_tv_s(t,S_t)}{B_t}.}
\]
Conversely, a sufficiently regular PDE solution with an admissible
self-financing strategy and the correct terminal payoff verifies the
replication. No arbitrage then identifies its initial cost as the claim
price. The argument uses one traded Brownian exposure to span the one
modeled Brownian risk.

\subsection{Gaussian solution, call formula, and Greeks}

Set \(\tau=T-t\). The logarithmic change of variable and removal of the
discounting term transform the PDE into a heat equation. Its Gaussian
representation is
\[
v(t,s)=e^{-r\tau}
\E\!\left[
g\!\left(s e^{(r-\sigma^2/2)\tau+\sigma\sqrt{\tau}Z}\right)
\right],
\qquad Z\sim N(0,1).
\]
For a call with strike \(K>0\), put
\[
d_1=\frac{\log(s/K)+(r+\sigma^2/2)\tau}{\sigma\sqrt{\tau}},
\qquad d_2=d_1-\sigma\sqrt{\tau}.
\]
The call, put, and parity formulas are
\[
\boxed{\begin{aligned}
C(t,s)&=s\Phi(d_1)-Ke^{-r\tau}\Phi(d_2),\\
P(t,s)&=Ke^{-r\tau}\Phi(-d_2)-s\Phi(-d_1),\\
C(t,s)-P(t,s)&=s-Ke^{-r\tau}.
\end{aligned}}
\]
Here \(\Phi\) and \(\varphi=\Phi'\) are the standard normal CDF and
density. The identity
\[
\boxed{s\varphi(d_1)=Ke^{-r\tau}\varphi(d_2)}
\]
cancels the density terms when differentiating. With calendar-time
\(\Theta=\partial_tC\), the call Greeks are
\[
\boxed{\begin{aligned}
\Delta_C&=\Phi(d_1),&
\Gamma_C&=\frac{\varphi(d_1)}{s\sigma\sqrt{\tau}},&
\mathrm{Vega}_C&=s\varphi(d_1)\sqrt{\tau},\\
\Theta_C&=-\frac{s\sigma\varphi(d_1)}{2\sqrt{\tau}}
           -rKe^{-r\tau}\Phi(d_2),&
\rho_C&=K\tau e^{-r\tau}\Phi(d_2).
\end{aligned}}
\]
Delta is the stock holding in the hedge. Gamma measures how quickly
that holding changes. Discrete rebalancing, funding frictions, jumps,
and volatility misspecification make the continuous-time hedge
imperfect.

\subsection*{Lecture 5 checklist}
\begin{itemize}[leftmargin=1.5em,itemsep=0.15em]
  \item distinguish portfolio value from self-financing gains;
  \item match diffusion and drift to obtain delta and the pricing PDE;
  \item recover the call formula and differentiate it to obtain the Greeks;
  \item check regularity, integrability, and the terminal payoff when
        verifying a candidate hedge.
\end{itemize}

% ============================================================
\section{Lecture 6 --- Reweighting, Girsanov, and Risk-Neutral Measure}
% ============================================================

\lecturesummary{A change of measure visits events with different
frequencies while an inverse density preserves the original
expectation. Gaussian tilting becomes Girsanov's drift shift for
Brownian paths. A true density martingale is needed before the
risk-neutral measure can be used.}

\subsection{Rare tails motivate a change of weights}

A change of measure can make a rare event easier to sample. To recover
its original probability, reweight the observations by the inverse density:
\[
\Pbb(A)=\E^{\Qbb}\!\left[
\frac{\dd\Pbb}{\dd\Qbb}\mathbf1_A\right]
\]
for equivalent measures \(\Pbb\) and \(\Qbb\).

\subsection{Densities and conditional Bayes' formula}

If \(Z_T>0\), \(\E^{\Pbb}[Z_T]=1\), and
\(\dd\Qbb/\dd\Pbb|_{\F_T}=Z_T\), then for integrable \(Y\),
\[
\boxed{\E^{\Qbb}[Y]=\E^{\Pbb}[Z_TY].}
\]
The density process is \(Z_t=\E^{\Pbb}[Z_T\mid\F_t]\). For
\(\mathcal G\subseteq\F_T\), the conditional version is
\[
\boxed{
\E^{\Qbb}[Y\mid\mathcal G]
=\frac{\E^{\Pbb}[Z_TY\mid\mathcal G]}
       {\E^{\Pbb}[Z_T\mid\mathcal G]}.}
\]
Changing measure changes the weights of paths, while equivalent
measures agree on null events.

\subsection{Stochastic exponential and the sign of Girsanov}

For a predictable process \(\lambda\) with finite pathwise energy, the
candidate density is
\[
\boxed{
Z_t=\exp\!\left(
-\int_0^t\lambda_u\dd W_u^{\Pbb}
-\tfrac12\int_0^t\lambda_u^2\dd u\right),
\qquad
\dd Z_t=-Z_t\lambda_t\dd W_t^{\Pbb}.}
\]
This is initially a nonnegative local martingale. A sufficient
condition for it to be a true martingale on \([0,T]\) is Novikov's
condition
\[
\E^{\Pbb}\!\left[
\exp\!\left(\tfrac12\int_0^T\lambda_u^2\dd u\right)\right]<\infty.
\]
If the density is a true martingale and
\(\dd\Qbb/\dd\Pbb|_{\F_T}=Z_T\), Girsanov gives
\[
\boxed{
W_t^{\Qbb}=W_t^{\Pbb}+\int_0^t\lambda_u\dd u.}
\]
Thus \(\dd W_t^{\Pbb}=\dd W_t^{\Qbb}-\lambda_t\dd t\). If
\(\dd X_t=b_t\dd t+a_t\dd W_t^{\Pbb}\), then
\[
\dd X_t=(b_t-a_t\lambda_t)\dd t+a_t\dd W_t^{\Qbb}.
\]
The drift changes; the quadratic variation
\(\dd[X]_t=a_t^2\dd t\) does not.

\subsection{The Black--Scholes pricing measure}

For the non-dividend-paying stock in Lecture~5, set
\(\lambda=(\mu-r)/\sigma\). With the corresponding true density
martingale,
\[
\boxed{
\frac{\dd S_t}{S_t}=r\dd t+\sigma\dd W_t^{\Qbb},
\qquad
\dd(e^{-rt}S_t)=\sigma e^{-rt}S_t\dd W_t^{\Qbb}.}
\]
The discounted price is a local martingale, and in this constant
coefficient model it is a true martingale. The measure \(\Qbb\)
is a pricing measure, not a forecast of historical returns. Its
conditional lognormal law also yields
\(C_t=S_t\Phi(d_1)-Ke^{-r(T-t)}\Phi(d_2)\) from two Gaussian
tails.

\subsection*{Lecture 6 checklist}
\begin{itemize}[leftmargin=1.5em,itemsep=0.15em]
  \item write a density and its inverse before reweighting an expectation;
  \item distinguish a local exponential from a true density martingale;
  \item keep the sign of the Brownian drift shift consistent with \(Z\);
  \item derive the risk-neutral drift of a non-dividend-paying stock.
\end{itemize}

% ============================================================
\section{Lecture 7 --- Martingale Prices, Hedges, and Completeness}
% ============================================================

\lecturesummary{Discounted self-financing gains lead to conditional
valuation when the relevant local martingale is a true martingale.
Dividends enter the stock gain and change its risk-neutral price drift.
Changing numeraire changes the units and measure together. Martingale
representation identifies a claim's Brownian exposure; traded assets
must span that exposure to implement a hedge.}

\subsection{From stochastic integrals to martingale prices}

Let \(B_t=\exp(\int_0^tr_u\dd u)\) and \(D_t=B_t^{-1}\).
For a non-dividend-paying stock under a money-market pricing measure,
\[
\frac{\dd S_t}{S_t}=r_t\dd t+\sigma_t\dd W_t^{\Qbb}
\quad\Longrightarrow\quad
\dd(D_tS_t)=D_t\sigma_tS_t\dd W_t^{\Qbb}.
\]
The discounted price is a local martingale because it is a stochastic
integral. For a self-financing portfolio
\(X_t=\Delta_tS_t+\beta_tB_t\), integration by parts gives
\[
\boxed{\dd(D_tX_t)=\Delta_t\dd(D_tS_t).}
\]
If this discounted wealth is a true martingale and \(X_T=H\), then
\[
\boxed{
V_t=X_t=B_t\E^{\Qbb}\!\left[\frac{H}{B_T}\middle|\F_t\right].}
\]
Local martingality alone does not justify the conditional-expectation
equality; integrability or an appropriate admissibility argument is
needed.

\subsection{Dividend gains and the pricing equation}

Suppose the stock pays continuous dividends \(qS_t\dd t\).
Its \emph{historical ex-dividend price} still has the form
\[
\dd S_t=\mu S_t\dd t+\sigma S_t\dd W_t^{\Pbb}.
\]
The gain from one share is \(\dd S_t+qS_t\dd t\). A
self-financing hedge therefore satisfies
\[
\boxed{
\dd X_t=\Delta_t(\dd S_t+qS_t\dd t)+\beta_t\dd B_t,}
\]
and its discounted form is
\[
\dd(D_tX_t)=\Delta_t\left[
\dd(D_tS_t)+D_tqS_t\dd t\right].
\]
For \(V_t=v(t,S_t)\), matching Brownian coefficients gives
\(\Delta_t=v_s\). Matching drifts cancels the historical price
drift \(\mu\) and gives
\[
\boxed{
v_t+(r-q)sv_s+\tfrac12\sigma^2s^2v_{ss}-rv=0.}
\]
The matching Brownian shift is
\(\lambda=(\mu+q-r)/\sigma\), so under the associated pricing measure
\[
\frac{\dd S_t}{S_t}=(r-q)\dd t+\sigma\dd W_t^{\Qbb}.
\]
The discounted \emph{gain} is driftless; the discounted
ex-dividend price alone has dividend-related drift.

\subsection{Change of numeraire}

For a strictly positive traded numeraire \(N\) with associated measure
\(\Qbb^N\), a replicable payoff has the representation
\[
\boxed{
V_t=N_t\E^{\Qbb^N}\!\left[\frac{H}{N_T}\middle|\F_t\right].}
\]
The numeraire, its measure, and the ratio \(H/N_T\) form one matched
choice of units. Taking \(N=B\) recovers money-market discounting.

\subsection{Martingale representation turns prices into hedges}

In a Brownian filtration, a square-integrable discounted claim value
has a representation
\[
M_t=\E^{\Qbb}[D_TH\mid\F_t]
=M_0+\int_0^t\psi_u^\top\dd W_u^{\Qbb}.
\]
The vector \(\psi_t\) is the claim's exposure to Brownian shocks. With
one non-dividend-paying traded stock satisfying
\(\dd(D_tS_t)=D_t\sigma_tS_t\dd W_t^{\Qbb}\), exposure matching gives
\[
\boxed{\Delta_t=\frac{\psi_t}{D_t\sigma_tS_t}.}
\]
If the claim is \(v(t,S_t)\), this reduces to \(\Delta_t=v_s(t,S_t)\).
For \(m\) assets driven by \(d\) Brownian factors, write
\(\dd\bar S_t=\Gamma_t\dd W_t^{\Qbb}\), with
\(\Gamma_t\in\R^{m\times d}\). Holdings \(\phi_t\) implement a
claim exposure \(\psi_t\) when
\[
\boxed{\Gamma_t^\top\phi_t=\psi_t.}
\]

\subsection{Completeness and the spanning condition}

In a \(d\)-factor Brownian market, all Brownian exposures can be spanned
when the discounted volatility matrix satisfies
\[
\boxed{\operatorname{rank}(\Gamma_t)=d
\quad\text{for almost every }(t,\omega).}
\]
Together with the required admissibility and integrability conditions,
this is the completeness criterion. A new traded asset helps only if
its Brownian exposure adds a missing direction.

When a Brownian direction is unspanned, its risk-premium shift can remain
free across equivalent martingale measures. Claims exposed to that
direction need not have a unique no-arbitrage price.

\subsection{Path dependence and the choice of state}

The full-filtration pricing identity applies to path-dependent claims.
To express the price as a function of the current state, include enough
variables to record the relevant path information. Apply the
multidimensional It\^o formula to this enlarged state: a finite-variation
state contributes a first derivative but no quadratic-variation or
covariation term.

\subsection*{Lecture 7 checklist}
\begin{itemize}[leftmargin=1.5em,itemsep=0.15em]
  \item distinguish discounted prices from discounted dividend gains;
  \item derive the dividend-adjusted PDE and delta hedge by replication;
  \item match a numeraire with its measure in a valuation formula;
  \item use martingale representation and a rank test to identify
        hedgeable Brownian exposures;
  \item identify the additional state variables needed for a path-dependent claim.
\end{itemize}

\paragraph{Examples to know.}
\begin{itemize}[leftmargin=1.5em,itemsep=0.15em]
  \item Price a unit payoff, a terminal share, a forward, and a digital;
        recover a European call by decomposing its payoff.
  \item Price and hedge a call with continuous dividends; recover
        put--call parity using a prepaid share.
  \item Choose a numeraire that simplifies the terminal payoff ratio.
  \item Find the martingale representations of linear and quadratic
        Brownian terminal payoffs.
  \item Analyze a stock driven by two Brownian risks and a claim depending
        on an untraded sector state; determine whether a second traded
        product completes the market.
  \item Use a nontraded claim to illustrate price ambiguity in an
        incomplete market.
  \item Introduce the running average for an arithmetic Asian option
        and derive its two-state pricing PDE.
\end{itemize}
